%! The Statistical Cycle and Types of Data — Unit 1, Lesson 1.1 — Student Packet Cover Sheet
% PILOT SHAPE (2026-09-06), Main Ideas / Notes table (2026-09-12): regenerated from the
% notes-only shape. THREE packet rows — Warm-Up, Guided Notes & Practice, Homework. The
% homework is the individual practice, started in class, and it is SCORED, so its cell is a
% \blank{}, never NA. The Guided Practice has no row of its own: it is the last row of the
% notes table and has no separate handout.
\documentclass[12pt]{article}
\usepackage{saar-article}
\usepackage{saar-boxes}
\usepackage{ltablex}
\keepXColumns
% Measured banner: the band grows to fit the title block, so a title long enough to wrap grows
% the band rather than running out of it (the stats course's \coverbanner; saar-article does not
% carry one, so it is defined here — every cover copies this block verbatim).
\newsavebox{\bannerbox}
\newlength{\bannerht}
\newlength{\coverbannerpad}
\setlength{\coverbannerpad}{0.16in}       % breathing room above and below the text
\newcommand{\coverbanner}[2]{%
  \sbox{\bannerbox}{%
    \begin{minipage}{\dimexpr\paperwidth-1.4in\relax}%
      \centering
      {\color{white}\textbf{\LARGE \CourseName}}\\[4pt]
      {\color{frostmid}\large\textbf{#1}}\\[6pt]
      {\color{white}\textbf{\large #2}}%
    \end{minipage}}%
  \setlength{\bannerht}{\dimexpr\ht\bannerbox+\dp\bannerbox+2\coverbannerpad\relax}%
  \noindent\begin{tikzpicture}[remember picture, overlay]
    \fill[cerulean] (current page.north west)
      rectangle ([yshift=-\bannerht]current page.north east);
    \node[anchor=north, inner sep=0pt]
      at ([yshift=-\coverbannerpad]current page.north) {\usebox{\bannerbox}};
  \end{tikzpicture}\par
  % The text body starts 0.6in below the page edge (geometry top=0.6in), so reserve only what
  % the band overruns that, plus a gap under the band.
  \vspace*{\dimexpr\bannerht-0.6in+0.16in\relax}%
}

\begin{document}

% Full-bleed cerulean banner, sized to the title block.
\coverbanner{Unit 1}{Lesson 1.1 \quad The Statistical Cycle and Types of Data}

\vspace{0.06in}
% NAMESTRIP: this is the ONLY component that carries the name row.
\namedateperiod
\vspace{0.18in}

% GRADUAL RELEASE: the vocabulary is taught in the guided notes, so the targets name it
% plainly. The standard codes (PS.DC.1a-c; AFDA.DA.2a) stay in the teacher-facing plan.
\begin{learningtargetbox}
\small
\begin{enumerate}[leftmargin=1.5em, itemsep=5pt]
  \item \ldots{} name the four stages of the \textbf{data cycle} and say what each stage
        \emph{hands to the next one}.
  \item \ldots{} test a question against three checks --- the answers \textbf{vary}, the
        \textbf{group} is named, the \textbf{measurement} is named --- and repair one that
        fails, naming the \textbf{variable of interest}.
  \item \ldots{} decide from the \emph{question alone}, before any data exist, whether the
        data will be \textbf{categorical} or \textbf{quantitative}, and build a
        \textbf{frequency table} with \textbf{relative frequencies}.
  \item \ldots{} write a \textbf{conclusion in context}, say what the data do \emph{not}
        allow anyone to claim, and name the \textbf{stage the problem entered at} when a
        study goes wrong.
\end{enumerate}
\end{learningtargetbox}

\vspace{0.14in}

\begin{tocbox}
\small
\renewcommand{\arraystretch}{1.5}
\begin{tabularx}{\linewidth}{@{} c l X r @{}}
\rowcolor{frost}
\textbf{\#} & \textbf{Component} & \textbf{Description} & \textbf{Score} \\
\midrule
1 & Warm-Up & The council's bus riders: the percent, the prediction, and two variables to
             classify & \blank{1.2cm} \\
2 & Guided Notes \& Practice & The four stages and what each hands on; the three tests a
             question must pass; the frequency table; why a correct percent can still fail to
             support a claim --- then the council's second survey, worked together
             & \blank{1.2cm} \\
% Row 3 is the lesson's GRADED practice, started in class. Packet pages are the default; a
% Desmos activity may replace them on a given day, decided at Close & Assign. The row and its
% score cell never change.
3 & Homework & Lakeside Farmers Market: name and classify the variable, organize and scale,
             and judge a report whose arithmetic is right ---
             \textbf{scored, started in class, due the first class after two study halls} & \blank{1.2cm} \\
\midrule
  & \multicolumn{2}{r}{\textbf{Total}} & \blank{1.2cm} \\
\end{tabularx}
\end{tocbox}

\vspace{0.14in}

% Keep in Mind carries the lesson's CONTENT — the rule, the test, the trap — never its process.
\begin{remindbox}
\small
The four stages run in order --- \textbf{formulate}, \textbf{collect}, \textbf{organize},
\textbf{analyze and communicate} --- and each stage can use only what the stage before it
handed over. \textbf{The question decides the type of data}, before anything is collected: if
what you record is a label, the variable is \textbf{categorical} and you summarize with
\textbf{relative frequencies}; if it is a number you can do arithmetic on, it is
\textbf{quantitative} and a mean is available. \textbf{Watch out:} a conclusion may answer
only the question that was asked, about only the group that was asked. When a study goes
wrong, name the stage the problem \emph{entered} at --- not stage 4, where the bad sentence
finally showed up.
\end{remindbox}

\end{document}
